% =====================================================================
\chapter[Maximum Likelihood and Maximum A Posteriori Estimation]{Maximum Likelihood and\\Maximum A Posteriori Estimation}
\markboth{Lab~\thechapter: MLE and MAP Estimation}{}
% =====================================================================

\labinfo{04\_MLE\_MAP}{mle\_map.py}{Diabetes progression (Efron et al., 2004): 442 patients, 10 baseline
measurements}{Regression with 64 polynomial features}{MAP: test RMSE 59.6, $R^2=0.385$; MLE: test RMSE
63.3, $R^2=0.308$}

\section{Dataset}
The \textbf{Diabetes} dataset (Efron, Hastie, Johnstone and Tibshirani, 2004) has 442 diabetes patients.
For each patient there are ten baseline measurements: age, sex, body mass index, average blood pressure
and six blood serum measurements (s1--s6). The target is a number (25--346) that measures how much the
disease progressed one year later.

To create a setting where MLE and MAP actually differ, I gave the model a lot of freedom: all ten
inputs, their squares (except for sex, which only has two values) and all 45 pairwise products, so
64 features in total. With only 310 training patients (70/30 split, 132 for testing), that is a lot of
weights to estimate from little data, which is exactly the situation where a prior should help.

\section{Code and experiments}
\begin{itemize}
\item \code{poly2} builds the 64 polynomial features, and \code{Prep} standardises them and centres $y$
      using training statistics only.
\item \code{fit\_mle} solves least squares (with \code{np.linalg.lstsq}, which also copes with the case
      of more weights than examples) and returns $\hat\sigma^2$.
\item \code{fit\_map} solves $(\Phi^\top\Phi+\lambda I)\vect w=\Phi^\top\vect y$.
\item \code{gaussian\_nll} measures how probable the test targets are under each fitted model.
\item \code{choose\_lambda} tries 50 values of $\lambda$ between $10^{-3}$ and $10^{4}$ with 5-fold
      cross-validation.
\end{itemize}
\noindent Besides the main comparison, the script draws a \emph{learning curve}: it trains both
estimators on random subsets of 60 to 310 patients (20 repeats each) and records the median test RMSE.

\section{Results}

Cross-validation chose $\lambda=7.20$. With the MAP noise estimate $\hat\sigma^2=2396$ this corresponds to
a prior standard deviation of $\tau=\sqrt{\sigma^2/\lambda}\approx18.3$ on each standardised weight.

\begin{table}[H]
\centering\small
\caption{MLE against MAP on the diabetes data (64 features, 310 training patients).}
\label{tab:mlemap}
\begin{tabular}{@{}l l S[table-format=2.2] S[table-format=2.2] S[table-format=1.4] S[table-format=1.3] S[table-format=4.1]@{}}
\toprule
\textbf{Estimator} & \textbf{Data} & {\textbf{MAE}} & {\textbf{RMSE}} & {\textbf{$R^2$}} & {\textbf{NLL}} & {$\|\vect w\|$}\\
\midrule
MLE & train & 34.99 & 44.98 & 0.6606 & 5.225 & 3209.7\\
MLE & test  & 51.38 & 63.25 & 0.3078 & 5.714 & 3209.7\\
\midrule
MAP & train & 38.97 & 48.95 & 0.5980 & 5.310 & 64.3\\
MAP & test  & \bfseries 49.36 & \bfseries 59.61 & \bfseries 0.3853 & \bfseries 5.551 & 64.3\\
\midrule
\multicolumn{7}{@{}l}{Reference: plain MLE with only the 10 raw inputs, test RMSE 58.66, $R^2=0.4046$}\\
\bottomrule
\end{tabular}
\end{table}

MLE fits the training patients better ($R^2$ 0.66 against 0.60) but generalises worse
(test $R^2$ 0.31 against 0.39). That is overfitting, plain and simple. The weight norms show why:
MLE uses weights fifty times larger than MAP, with big positive and negative values that cancel each
other on the training data but not on new patients. MAP also gives the test targets a higher
probability (lower negative log-likelihood).

This is the classic regularisation picture. Training error only goes up as
$\lambda$ grows, but test and CV error first go down (less overfitting) and then up again (the prior is
so strong that the model cannot fit anything).

\begin{table}[H]
\centering\small
\caption{Median test RMSE for small training sets (20 random subsets each).}
\label{tab:mlemap-lc}
\begin{tabular}{@{}l c c c c c c@{}}
\toprule
Training patients & 60 & 80 & 100 & 150 & 200 & 310\\
\midrule
MLE & 543.5 & 190.7 & 118.3 & 77.3 & 67.1 & 63.3\\
MAP & 59.8 & 60.0 & 59.1 & 59.6 & 59.4 & 59.6\\
\bottomrule
\end{tabular}
\end{table}

The learning curve shows the biggest difference. With 60 patients there are fewer equations than
weights (64), MLE has infinitely many perfect fits, and the one it returns is useless (RMSE 543 on a
target whose standard deviation is 77). MAP is almost unaffected by the amount of data; the prior
supplies what the data cannot. As the training set grows, MLE slowly catches up, as theory says it should.

One honest observation: the plain model with only the ten raw inputs (test $R^2=0.40$) is still slightly
better than MAP with 64 features. The extra polynomial features do not carry much new information for this
dataset, and MAP's main achievement is to stop them from doing harm.

\begin{figure}[H]
\centering
\includegraphics[width=0.48\linewidth]{\figpath{04\_MLE\_MAP}{rmse_vs_lambda.pdf}}\hfill
\includegraphics[width=0.48\linewidth]{\figpath{04\_MLE\_MAP}{learning_curve.pdf}}
\caption{The effect of prior strength and training-set size on MLE and MAP.}
\end{figure}

\section{Source code}
\lstinputlisting[caption={\texttt{04\_MLE\_MAP/mle\_map.py}}]{04_MLE_MAP/mle_map.py}
\clearpage
